% TODO make hw stylesheet
\documentclass[10pt]{article}
\usepackage[letterpaper,top=1in,left=1in,right=1in]{geometry}
\usepackage[dvipsnames,svgnames]{xcolor}
\usepackage[shortlabels]{enumitem}
\usepackage[framemethod=TikZ]{mdframed}
\usepackage{amsmath,amssymb,amsthm}
\usepackage[colorlinks]{hyperref}
\usepackage{multicol}
\usepackage{tabularx}
\usepackage{bbm}

\hypersetup{
    colorlinks=true,
    linkcolor=blue,
    filecolor=magenta,
    urlcolor=magenta,
    pdftitle={MAT 549 HW}
}

\usepackage{mathtools}
\usepackage{thmtools}
\usepackage[textsize=footnotesize]{todonotes}
\usepackage{titling}
\usepackage{cancel}

\reversemarginpar
\mdfdefinestyle{mdbluebox}{roundcorner=10pt,innerbottommargin=9pt,
linecolor=blue,backgroundcolor=TealBlue!5,}
\declaretheoremstyle[headfont=\sffamily\bfseries\color{MidnightBlue},
mdframed={style=mdbluebox},]{thmbluebox}

\mdfdefinestyle{mdbloobox}{frametitlefont=\bfseries,innerbottommargin=8pt,
nobreak=true,backgroundcolor=TealBlue!5,linecolor=blue,}
\declaretheoremstyle[headfont=\bfseries\color{MidnightBlue},
mdframed={style=mdbloobox},headpunct={\\[3pt]},postheadspace=0pt,]{thmbloobox}

\mdfdefinestyle{mdredbox}{frametitlefont=\bfseries,innerbottommargin=8pt,
nobreak=true,backgroundcolor=Salmon!5,linecolor=RawSienna,}
\declaretheoremstyle[headfont=\bfseries\color{RawSienna},
mdframed={style=mdredbox},headpunct={\\[3pt]},postheadspace=0pt,]{thmredbox}

\mdfdefinestyle{mdgreenbox}{linecolor=ForestGreen,backgroundcolor=ForestGreen!5,
linewidth=2pt,rightline=false,leftline=true,topline=false,bottomline=false,}
\declaretheoremstyle[headfont=\bfseries\sffamily\color{ForestGreen!70!black},
mdframed={style=mdgreenbox},headpunct={ --- },]{thmgreenbox}

\mdfdefinestyle{mdorangebox}{linecolor=RedOrange,backgroundcolor=RedOrange!5,
linewidth=2pt,rightline=false,leftline=true,topline=false,bottomline=false,}
\declaretheoremstyle[headfont=\bfseries\sffamily\color{RedOrange!70!black},
mdframed={style=mdorangebox},headpunct={ --- },]{thmorangebox}

\mdfdefinestyle{mdblackbox}{linecolor=black,backgroundcolor=RedViolet!5!gray!5,
linewidth=3pt,rightline=false,leftline=true,topline=false,bottomline=false,}
\declaretheoremstyle[mdframed={style=mdblackbox}]{thmblackbox}

\declaretheoremstyle[spaceabove=6pt,spacebelow=6pt]{boldhead}
\declaretheoremstyle[spaceabove=6pt,spacebelow=6pt,headfont=\normalfont\itshape]{italichead}

\declaretheorem[style=boldhead,name=Problem]{problem}
\declaretheorem[style=boldhead,name=Problem,numbered=no]{problem*}
\declaretheorem[style=boldhead,name=Setup]{setup} % for config geo
\declaretheorem[style=boldhead,name=Example,sibling=problem]{example}
\declaretheorem[style=boldhead,name=$\bigstar$ Problem,sibling=problem]{niceprob}
\declaretheorem[style=boldhead,name=Theorem,sibling=problem]{theorem}
\declaretheorem[style=boldhead,name=Corollary,sibling=theorem]{corollary}
\declaretheorem[style=boldhead,name=Proposition,sibling=theorem]{prop}
\declaretheorem[style=boldhead,name=Lemma,numberwithin=problem]{lemma}
\declaretheorem[style=boldhead,name=Theorem,numbered=no]{theorem*}
\declaretheorem[style=boldhead,name=Lemma,numbered=no]{lemma*}
\declaretheorem[style=boldhead,name=Claim,numberwithin=problem]{claim}
\declaretheorem[style=boldhead,name=Claim,numbered=no]{claim*}
\declaretheorem[style=boldhead,name=Remark,numberwithin=problem]{remark}
\declaretheorem[style=boldhead,name=Remark,numbered=no]{remark*}
\declaretheorem[style=boldhead,name=Definition,sibling=theorem]{definition}
\declaretheorem[style=boldhead,name=Definition,numbered=no]{definition*}

\providecommand{\ol}{\overline}
\providecommand{\ceil}[1]{\left\lceil #1 \right\rceil}
\providecommand{\floor}[1]{\left\lfloor #1 \right\rfloor}
\newcommand{\dd}{\mathrm{d}}
\providecommand{\ul}{\underline}
\providecommand{\eps}{\varepsilon}
\providecommand{\half}{\frac{1}{2}}
\providecommand{\CC}{\mathbb C}
\providecommand{\FF}{\mathbb F}
\providecommand{\II}{\mathbb I}
\providecommand{\NN}{\mathbb N}
\providecommand{\QQ}{\mathbb Q}
\providecommand{\RR}{\mathbb R}
\providecommand{\ZZ}{\mathbb Z}
\providecommand{\ind}{\mathbbm 1}
\providecommand{\dg}{^\circ}
\providecommand{\ii}{\item}
\providecommand{\setdiff}{\smallsetminus}
\providecommand{\alert}[1]{{\sffamily\textbf{\textcolor{blue}{#1}}}}
\providecommand{\inv}{^{-1}}
\providecommand{\mb}[1]{\mathbf{#1}}
\newcommand{\what}{\widehat}

\newenvironment{soln}{\begin{proof}[Solution]}{\end{proof}}

\providecommand{\pow}{\operatorname{Pow}}
\providecommand{\vocab}[1]{{\textbf{\textcolor{ForestGreen}{#1}}}}
\providecommand{\lcm}{\operatorname{lcm}}
\providecommand{\abs}[1]{\left\lvert #1\right\rvert}
\providecommand{\smabs}[1]{\lvert #1\rvert}
\providecommand{\innerprod}[1]{\langle\!\langle #1\rangle\!\rangle}
\providecommand{\norm}[1]{\left\| #1\right\|}
\providecommand{\smnorm}[1]{\| #1\|}
\providecommand{\gr}{\operatorname{gr}}

\usepackage{mathrsfs}
\providecommand{\tanvec}[1]{\frac{\partial}{\partial #1}}

% place title at top right
\setlength{\droptitle}{-8ex}
\pretitle{\begin{flushright}\Large\bfseries}
\posttitle{\par\end{flushright}}
\preauthor{\begin{flushright}}
\postauthor{\end{flushright}}
\predate{\begin{flushright}}
\postdate{\end{flushright}}

\title{MAT 549 HW03}
\author{\large Aditya Pahuja}
\date{\large Due 03/03}
\begin{document}
\maketitle

\begin{problem}
    [10-6, vector bundle construction theorem]
    Let $M$ be a smooth manifold and let $\{U_\alpha\}_{\alpha\in A}$
    be an open cover of $M$. Suppose that for each
    $\alpha,\beta\in A$ we are given a smooth map
    $\tau_{\alpha\beta}\colon U_{\alpha}\cap U_\beta\to GL_k\RR$
    such that $\tau_{\alpha\beta}\tau_{\beta\gamma} = \tau_{\alpha\gamma}$
    on $U_\alpha\cap U_\beta\cap U_\gamma$ for all
    $\alpha,\beta,\gamma\in A$. Show that there is
    a smooth rank-$k$ vector bundle $p\colon E\to M$
    with smooth local trivializations
    $\Phi_\alpha\colon p\inv(U_\alpha)\to U_\alpha\times\RR^k$
    whose transition functions are the given maps $\tau_{\alpha\beta}$.
\end{problem}

\begin{soln}
    Let $E = (\coprod_{\alpha\in A} U_\alpha\times\RR^k)/{\sim}$
    where $(x_1,v_1)\in U_\alpha\times\RR_k$
    and $(x_2,v_2)\in U_\beta\times\RR^k$
    satisfy $(x_1,v_1)\sim(x_2,v_2)$ if
    $(x_1,v_1) = (x_2,\tau_{\alpha\beta}(x_1)v_2)$
    (in particular $x_1$ must equal $x_2$).
    Observe that $\sim$ is an equivalence relation:
    \begin{itemize}
        \ii Let $(x,v)\in U_\alpha\times\RR^k$.
        Note that
        $\tau_{\alpha\alpha}(x) = \operatorname{id}_{x_\alpha\times\RR^k}$
        since $\tau_{\alpha\alpha}(x)^2 = \tau_{\alpha\alpha}(x)$
        for each $x\in U_\alpha$ by the cocycle condition and
        the invertibility of $\tau_{\alpha\alpha}(x)\in GL_k\RR$,
        so $(x,v) = (x,\tau_{\alpha\alpha}(x)v)$
        Thus $(x,v)\sim(x,v)$.
        \ii Let $(x,v_1),(x,v_2)\in (U_\alpha\cap U_\beta)\times\RR^k$
        such that $(x,v_1)\sim(x,v_2)$. Then observe that
        \begin{equation*}
            \tau_{\alpha\beta}(x)\tau_{\beta\alpha}(x) = \tau_{\alpha\alpha}(x)
            = \operatorname{id}_{x\times\RR^k},\qquad
            \tau_{\beta\alpha}(x)\tau_{\alpha\beta}(x) = \tau_{\beta\beta}(x)
            = \operatorname{id}_{x\times\RR^k}
        \end{equation*}
        for $x\in U_\alpha\cap U_\beta$, so $\tau_{\alpha\beta}(x)$
        and $\tau_{\beta\alpha}(x)$ are inverse linear transformations. Thus
        \begin{equation*}
            v_1 = \tau_{\alpha\beta}(x)v_2
            \iff v_2 = \tau_{\beta\alpha}(x)v_1
        \end{equation*}
        which means that $(x,v_2)\sim(x,v_1)$.
        \ii Let $(x,v_1),(x,v_2),(x,v_3)\in
        (U_\alpha\cap U_\beta\cap U_\gamma) \times\RR^k$,
        such that $(x,v_1)\sim (x,v_2)$ and $(x,v_2)\sim(x,v_3)$. Then
        \begin{equation*}
            v_1 = \tau_{\alpha\beta}(x)v_2
            = \tau_{\alpha\beta}(x)\tau_{\beta\gamma}(x)v_3
            = \tau_{\alpha\gamma}(x)v_3
        \end{equation*}
        by the cocycle condition, so $(x,v_1)\sim(x,v_3)$.
    \end{itemize}

    Define $p\colon E\to M$ by sending the equivalence class of
    a pair $(x,v)$ to $x$; this is well-defined because
    two equivalent elements in $\coprod_{\alpha\in A}U_\alpha\times\RR^k$
    must have the same first coordinate by definition.
    Then, define
    \begin{equation*}
        \Phi_\alpha\colon p\inv(U_\alpha)\to U_\alpha\times\RR^k,\qquad
	\Phi_\alpha([(x,v)]) = (x,v)
    \end{equation*}
    where $[(x,v)]$ is the equivalence class of
    $(x,v)\in U_\alpha\times\RR^k$;
    this is well-defined because if $(x,v_1)\sim(x,v_2)$
    for $(x,v_1),(x,v_2)\in U_\alpha\times\RR^k$, then
    $v_1 = \tau_{\alpha\alpha}(x)v_2 = v_2$,
    i.e. each equivalence class in $E$ contains at most one element
    of each set in the disjoint union, and it is a bijection
    because it has an inverse given by just sending $(x,v)$
    to the equivalence class containing it. We can then see that
    for $(x,v)\in (U_\alpha\cap U_\beta)\times\RR^k$,
    \begin{equation*}
        \Phi_\alpha\circ\Phi_\beta\inv(x,v)
	= \Phi_\alpha([(x,v)])
	= \Phi_\alpha([(x,\tau_{\alpha\beta}(x)v)])
	= (x,\tau_{\alpha\beta}(x)v)
    \end{equation*}
    where $(x,\tau_{\alpha\beta}(x)v)$ is the unique element
    of $U_\alpha\times\RR^k$ which is equivalent to
    $(x,v)\in U_\beta\times\RR^k$ under $\sim$
    (uniqueness follows from the previous sentence,
    since the equivalence class contains at most
    one element of $U_\alpha\times\RR^k$).

    The work above shows that $E$, together with the given data,
    satisfies the conditions of Lee's Lemma 10.6,
    so $E$ has a unique topology and smooth structure
    making it into a smooth rank-$k$ vector bundle
    whose local trivializations are the $\Phi_\alpha$;
    and of course the transition function
    between $\Phi_\alpha$ and $\Phi_\beta$ is $\tau_{\alpha\beta}$
    by definition of transition function.
\end{soln}

\begin{problem}
    [10-7]
    Compute the transition function for $TS^2$ associated with
    the trivializations determined by stereographic coordinates.
\end{problem}

\begin{soln}
    Let $\sigma$ be the stereographic projection from the north pole
    and $\tilde\sigma = -\sigma(-x)$ the stereographic projection
    from the south pole, as per Lee's Problem 1-7.
    Then, we can compute
    \begin{align*}
        \tilde\sigma(\sigma\inv(u_1,u_2))
	&= \tilde\sigma\left(\frac{u_1}{\abs u^2+1},
	\frac{u_2}{\abs u^2+1},\frac{\abs u^2-1}{\abs u^2+1}\right)\\
	&= -\frac 1{1 + \frac{\abs u^2-1}{\abs u^2+1}}
	\left(-\frac{u_1}{\abs u^2+1},-\frac{u_2}{\abs u^2+1}\right)
	= \frac 1{\abs u^2}(u_1,u_2)
    \end{align*}
    as the change of coordinates between $S^2$ minus the north pole $N$
    and $S^2$ minus the south pole $S$,
    so that at each point $p$ of $S^2 - \{N\} - \{S\}$,
    the transition function is simply the map between tangent spaces
    induced by $\tilde\sigma\circ\sigma\inv$, namely the Jacobian
    \begin{equation*}
	\frac 1{(u_1^2+u_2^2)^2}
        \begin{pmatrix}
	    u_2^2-u_1^2&-2u_1u_2\\-2u_1u_2&u_1^2-u_2^2
        \end{pmatrix}.\qedhere
    \end{equation*}
\end{soln}

\begin{problem}
    [11-7]
    In the following problems, $M$ and $N$ are smooth manifolds,
    $F\colon M\to N$ is a smooth map, and $\omega\in\mathfrak X^*(N)$.
    Compute $F^*\omega$ in each case.
    \begin{enumerate}[(a)]
        \ii $M = N = \RR^2$; $F(s,t) = (st,e^t)$;
	$\omega = x\dd y - y\dd x$.
	\ii $M = \RR^2$; $N = \RR^3$;
	$F(\theta,\varphi) = ((\cos\varphi+2)\cos\theta,
	(\cos\varphi+2)\sin\theta, \sin\varphi)$;
	$\omega = z^2\dd x$.
	\ii $M = \{(s,t)\in\RR^2 : s^2 + t^2 < 1\}$;
	$N = \RR^3 - \{0\}$;
	$F(s,t) = (s,t,\sqrt{1-s^2-t^2})$;
	$\omega = (1-x^2-y^2)\dd z$.
    \end{enumerate}
\end{problem}

\begin{soln}
    For each $p\in M$ and $v\in T_pM$,
    \begin{equation*}
	(F^*\omega)_p(v) = \omega_{F(p)}(\dd F_p(v)).
    \end{equation*}
    (a) Since
    \begin{equation*}
	\dd F_{(s,t)} = \begin{pmatrix}
	    t&s\\
	    0&e^t
        \end{pmatrix}
    \end{equation*}
    we can compute
    \begin{equation*}
	(F^*\omega)_{(s,t)}\begin{pmatrix}
	    a\\b
	\end{pmatrix}
	= \omega_{F(s,t)} \begin{pmatrix}
	    at+bs\\ be^t
	\end{pmatrix}
	= stbe^t - e^t(at+bs) = -te^t\cdot a + se^t(t-1)\cdot b,
    \end{equation*}
    so $F^*\omega = -te^t\dd s + se^t(t-1)\dd t$.\\

    (b) Since
    \begin{equation*}
	\dd F_{(\theta,\varphi)} = \begin{pmatrix}
	    -(\cos\varphi+2)\sin\theta&-\sin\varphi\cos\theta\\
	    (\cos\varphi+2)\cos\theta&-\sin\varphi\sin\theta\\
	    0 & \cos\varphi
	\end{pmatrix}
    \end{equation*}
    we can compute
    \begin{equation*}
	(F^*\omega)_{(\theta,\varphi)}\begin{pmatrix}
	    a\\b
	\end{pmatrix}
	= \omega_{F(\theta,\varphi)}\begin{pmatrix}
	    -a(\cos\varphi+2)\sin\theta - b\sin\varphi\cos\theta\\
	    a(\cos\varphi+2)\cos\theta - b\sin\varphi\sin\theta\\
	    b\cos\varphi
	\end{pmatrix}
	= -\sin^2\varphi(\cos\varphi+2)\sin\theta\cdot a
	- \sin^3\varphi\cos\theta\cdot b
    \end{equation*}
    so $F^*\omega = -\sin^2\varphi(\cos\varphi+2)\sin\theta\dd\theta
    - \sin^3\cos\theta\dd\varphi$.
    \\

    (c) Since
    \begin{equation*}
	\dd F_{(s,t)} = \begin{pmatrix}
	    1&0\\0&1\\
	    \frac{-s}{\sqrt{1-s^2-t^2}}&
	    \frac{-t}{\sqrt{1-s^2-t^2}}
	\end{pmatrix}
    \end{equation*}
    we can compute
    \begin{equation*}
	(F^*\omega)_{(s,t)}\begin{pmatrix}
	    a\\b
	\end{pmatrix}
	= \omega_{F(s,t)}\begin{pmatrix}
	    a\\b\\\frac{-as-bt}{\sqrt{1-s^2-t^2}}
	\end{pmatrix}
	= -s\sqrt{1-s^2-t^2}\cdot a - t\sqrt{1-s^2-t^2}\cdot b
    \end{equation*}
    so $F^*\omega = -s\sqrt{1-s^2-t^2}\dd s
    - t\sqrt{1-s^2-t^2}\dd t$.
\end{soln}

\begin{problem}
    [11-9]
    Let $f\colon\RR^3\to\RR$ be the function
    $f(x,y,z) = x^2 + y^2 + z^2$, and let
    $F\colon\RR^2\to\RR^3$ be
    \begin{equation*}
	F(u,v) = \left(\frac{2u}{u^2+v^2+1}, \frac{2v}{u^2+v^2+1},
	\frac{u^2+v^2-1}{u^2+v^2+1}\right),
    \end{equation*}
    the inverse of the stereographic projection.
    compute $F^*\dd f$ and $\dd(f\circ F)$ separately,
    and verify that they are equal.
\end{problem}

\begin{soln}
    Noting that $\dd f = 2x\dd x + 2y\dd y + 2z\dd z$, we can calculate
    \begin{align*}
        F^*(2x\dd x)
	&= \frac{2u}{u^2+v^2+1}\left(\frac{2(-u^2+v^2+1)}{(u^2+v^2+1)^2}\dd u
	- \frac{4uv}{(u^2+v^2+1)^2}\dd v\right)\\
        F^*(2y\dd y)
	&= \frac{2v}{u^2+v^2+1}\left(- \frac{4uv}{(u^2+v^2+1)^2}\dd u
	+ \frac{2(u^2-v^2+1)}{(u^2+v^2+1)^2}\dd v\right)\\
        F^*(2z\dd z)
	&= \frac{u^2+v^2-1}{u^2+v^2+1}\left(\frac{4u}{(u^2+v^2+1)^2}\dd u
	+ \frac{4u}{(u^2+v^2+1)^2}\dd v\right).
    \end{align*}
    Summing, we find that the coefficient of $\dd u$ is
    $(u^2+v^2+1)^{-3}$ times
    \begin{equation*}
        4u(-u^2+v^2+1) - 8uv^2 + 4u(u^2+v^2-1)
	= 4u(2v^2) - 8uv^2 = 0,
    \end{equation*}
    and by symmetry the coefficient of $\dd v$ is also zero,
    so $F^*\dd f \equiv 0$.

    On the other hand, $f\circ F = 1$
    since $F$ sends points in $\RR^2$ to the sphere in $\RR^3$
    by definition of (the inverse of) stereographic projection,
    implying $\dd(f\circ F) = 0$, as expected from above.
\end{soln}











\end{document}
